\subsection{Proof of Injectivity}
\label{subsections:techniques:universal:injective}



\paragraph{Description.}
A proof of injectivity proves that a function is one-to-one, %
meaning that no two different inputs have the same output. %
Precisely, it proves that a function $f : X \to Y$ satisfies %
$\forall x_1, x_2 \in X$, if $f(x_1) = f(x_2)$, then $x_1 = x_2$. %
Generally, a proof of injectivity mirrors the following format:
\begin{enumerate}{
    \item \textbf{Setup.} %
        Let $x_1, x_2 \in X$ such that $f(x_1) = f(x_2)$.
    \item \textbf{Verification.} %
        Show that $x_1 = x_2$.
    \item \textbf{Conclusion.} %
        Conclude the proof, optionally by summarizing what has been found.
}
\end{enumerate}



\paragraph{Advice.}
\begin{itemize}{
    \item Use this proof for a statement that a function is one-to-one, %
        also called injective.
    \item This proof is largely algebra and simplification: %
        simplify $f(x_1) = f(x_2)$ until $x_1 = x_2$ is reached.
}
\end{itemize}



\paragraph{Example.}
\begin{exercise}{
    Let $f : \Z \to \Z$ be defined by $f(z) = 2z + 1$. %
    Prove that $f$ is one-to-one.
}
\end{exercise}
\begin{proof}{
    This statement can be proven by showing that $f$ is one-to-one.

    \begin{annotatedpanel}[Setup.]{
        Let $z_1, z_2 \in \Z$ such that $f(z_1) = f(z_2)$.
    }{
        \textbf{Step 1.} Let $x_1, x_2 \in X$ such that %
        $f(x_1) = f(x_2)$. %
        Here, $X = \Z$, $x_1 = z_1$, and $x_2 = z_2$.
    }
    \end{annotatedpanel}

    \begin{annotatedpanel}[Verification.]{
        Then $2z_1 + 1 = 2z_2 + 1$. %
        Thus, $2z_1 = 2z_2$, so $z_1 = z_2$.
    }{
        \textbf{Step 2.} Show that $x_1 = x_2$. %
        Here, subtracting $1$ and dividing by $2$ gives %
        $z_1 = z_2$.
    }
    \end{annotatedpanel}

    \begin{annotatedpanel}[Conclusion.]{
        Thus, $\forall z_1, z_2 \in \Z$, if $f(z_1) = f(z_2)$, %
        then $z_1 = z_2$. %
        Therefore, the statement holds by the definition of one-to-one. %
        \qedhere
    }{
        \textbf{Step 3.} Conclude the proof. %
        Here, the summary states what has been found: %
        equal outputs have equal inputs.
    }
    \end{annotatedpanel}
}
\end{proof}